\section{Interchanging address fields on tapes}
\label{sec:interchange}

A square matrix stored row by row becomes its transpose when the row
and column fields of its address exchange places. Multiplication will
need the same operation on selected fields of a much longer address.
For example, the entry at $(i,x,g,y,t)$ must be written at
$(i,y,g,x,t)$. The blocks indexed by $i$, the intervening index $g$,
and the records indexed by $t$ can all be arbitrarily large. Their
total volume affects the time linearly. It is the width of $x$ and
$y$ that determines the overhead per symbol.

Write $[d]=\{0,\ldots,d-1\}$ and store arrays in lexicographic order,
most significant index first. The address is the position of a symbol
in a stream. The machine never needs a separate address record for
each symbol.

\begin{theorem}[Interchange on fixed tapes]\label{thm:interchange}
For each $\alpha>1/2$ there is a deterministic multitape Turing machine
with a fixed finite alphabet and a fixed number of one-dimensional tapes
with the following property. Given positive integers $P,G,B,u$ in
canonical binary and a bit array in lexicographic order on
\[
          [P]\times[2^u]\times[G]\times[2^u]\times[B],
\]
it writes the array obtained by interchanging the second and fourth
indices. If $V=PGB\,2^{2u}$, it uses $O_\alpha(Vu^\alpha)$ transitions
and visits $O_\alpha(V)$ tape cells. In each transition each tape head
moves by at most one cell. The same bounds hold for a stream over any
fixed finite alphabet, with constants also depending on that alphabet.
\end{theorem}

For $P=G=B=1$, Theorem~\ref{thm:interchange} transposes a binary
matrix of side $2^u$. A digit-by-digit implementation of the swap
makes a separate linear-volume pass for each of the $u$ pairs of
digits. Our algorithm shares a smaller collection of passes among
many pairs. We prove the result for $1/2<\alpha<1$, which implies all
larger exponents as well.

\subsection{Tape operations and the target for one recursive step}

The recursion uses a prime radix $b$ chosen below. A field of $f$
consecutive digits takes values in $\mathbb Z/b^f\mathbb Z$. A
\emph{call} interchanges two such fields in a single stream, leaving
all other coordinates fixed. At each level we charge separately for
calls and for operations that take linear time in the stream volume.

Copying aligned streams, combining their entries by a fixed finite
table, and splitting and merging streams at a fixed digit all take
linear time. More generally, an invertible lower triangular rational
change on a fixed number of fields, read in physical order, costs a
linear number of tape steps. Apply the changes in descending order of
the target field, so each controlling field still has its original
value. A diagonal scaling acts on one field and each off-diagonal
term shifts it by a multiple of an earlier field. These operations
include preparation of the shifts in the presence of arbitrary
intervening and trailing indices \cite[Lemmas~3 and~4]{Tape}.

For instance, a fixed value of the controlling field precedes a run of
at least $b^f$ entries containing the shifted field. This run pays for
reading the control and computing the wrap point. Rotating contiguous
blocks implements the shift. Multiplication by a fixed positive
integer that is a unit modulo $b$ has a bounded number of increasing
pieces, and a finite residue counter enumerates the outputs. Reversing
the construction gives inverse scaling. For nested scans a binary
counter on a separate tape tracks each range. Ripple carries, resets,
and head returns cost linear time over each complete scan
\cite[Section~2]{Tape}. For a general invertible rational change on
$d$ fields, factor its matrix as $L_1\Pi L_2$ with lower triangular
$L_1,L_2$ and a field permutation $\Pi$. This requires at most $d-1$
calls in addition to the linear work \cite[Lemma~9]{Tape}. For this factorization see also
\cite{Grigoriev1982}. We choose $b$ after all the rational matrices and
factorizations, avoiding their denominators and the numerators of their
required units. They then specialize to every modulus $b^f$.

At a recursive step the two chunks each contain $m$ main fields of
width $f$. Let $Q_j$ exchange the $j$th pair. Their product
interchanges all main fields. On a fixed batch of independent streams,
the prediction and exact repair of
\cite[Theorem~\ext{F}{thm:repair}]{oaF} reduces this product to
$O(\sqrt{m\log(2m)})$ stages, together with conversions on a small
fraction of the streams. Each stage combines the identity and the
individual swaps linearly. Executing its $m$ swaps separately would
use up the saving in the number of stages.

To implement one stage, assign a label to each batch. For the $j$th
swap each source and destination label will have its own way of
expressing its physical tape address in a common coordinate system.
Sending a stream to its own label in those coordinates performs the
swap. We implement the identity on labels as the difference of two
maps. The first mixes the labels through a matrix of small rank. The
second cancels the unwanted off-diagonal transfers, whose support is
partitioned into rectangles. In a rectangle every source sends to
every destination, so we can add the sources together before making
any copies. The cancellation is valid only if both contributions
refer to the same logical address.

Two quantities therefore matter for the choice of labels. The
rank \emph{fraction} controls how many streams enter the full
coordinate changes for individual swaps. The number of rectangles
incident to one label and the dimension of an address change control
the work shared by all swaps. For every $\eta>0$ we can arrange that
the fraction is $O(1/m)$ and that the latter two quantities are
$O(m^\eta)$ each. The calls per stream in one stage will then be
$O(m^{2\eta})$. Even a very large number of labels costs only a fixed factor in the
scans, since the labels are chosen before the recursion starts. This use of different fields for
label transfers and address changes builds on \cite{Tape}. For the
index-coding interpretation see
\cite{BarYossefBirkJayramKol2006,Lubetzky}, and for the intersection
constructions underlying the two-field method see
\cite{FranklWilson1981,Alon1998,Matousek2010} and
\cite[Section~1.6]{Tape}.

\subsection{A family of labels}

Choose an odd prime $p$ and a power of two $q$. Here $p$ is the
characteristic of the data. The address maps will be built over
$\mathbb Q$. Set $h=2p-1$ and let $\mathcal U=[q]^h$, of size
$n_0=q^h$. Call two elements of $\mathcal U$ far if their Hamming
distance is $p$. A label is an element of
$\mathcal V=\mathcal U^h$, of size $N=n_0^h$. For $v,w\in\mathcal V$
let $k(v,w)$ be the number of coordinates in which the two inner words
are far. Define the symmetric matrix $P_*$ on $\mathcal V$ by assigning
a one to each diagonal entry and to each pair differing in exactly $p$
coordinates with far inner words in all $p$. Give the other entries
value zero.

\begin{lemma}\label{lem:interchange-labels}
The rank $r_*$ of $P_*$ over $\mathbb F_p$ satisfies
$r_*/N=O_p(q^{-p^2})$. Its off-diagonal support is a disjoint union
of directed rectangles. Each label belongs to at most
\[
          \lambda=O_p\big(q^p(1+\log q)^{p^2}\big)
\]
rectangles on either side. On each rectangle $k(v,w)=p$ across the
two sides, and $k(v,w)=0$ within either side.
\end{lemma}

\begin{proof}
For inner words at distance $d$, the matrix $P_0$ of equal or far
pairs has entry $1-d^{p-1}$ in $\mathbb F_p$. Write $d$ as the sum
of coordinate inequality indicators and expand. Each resulting term
depends on at most $p-1$ coordinates. It factors as a tensor product
of matrices of rank at most $q$ on those coordinates and all-ones
matrices on the others. There are $O_p(1)$ terms, so
$\operatorname{rank}(P_0)/n_0=O_p(q^{-p})$.
Writing $e$ for the outer distance, $P_*$ is the entrywise product of
$P_0^{\otimes h}$ and the matrix with entry $1-e^{p-1}$. Expand the
latter in outer coordinate indicators. In each term at least
$h-(p-1)=p$ coordinates retain $P_0$, and each other coordinate has
rank at most $n_0$. The tensor product rank bound and summation over
$O_p(1)$ terms give the claimed estimate.

For an ordered far pair of inner words, specify the $p$ differing
positions and the common letters elsewhere. At the first differing
position fix both letters. At each other differing position partition
ordered pairs of distinct letters into rectangles using the first bit
at which their binary strings differ. A rectangle specifies the common
prefix and the orientation of the two next bits. Each letter belongs
to at most $\log_2 q$ such rectangles on a specified side. This gives
a disjoint inner decomposition with
$O_p(q(1+\log q)^{p-1})$ incidences per word on either side. Words
on one side can vary in only $p-1$ positions, so no two are far.
For the outer pairs specify the $p$ differing coordinates, fix the
common inner words elsewhere, and take a product of these inner
rectangles on the differing coordinates. This proves the incidence
bound and both assertions about $k$.
\end{proof}

The scalar transfers will take place over $\mathbb F_p$. The same
labels must also describe coordinate changes which can be prepared on
tapes. For this purpose we use a rational vector space in which two
labels on the same side of a rectangle have the same pairwise Gram
block, while going to the other side exchanges the two coordinates.

On the rational vector space with two generators per label take the
symmetric bilinear form whose $(v,w)$ block is
\begin{equation}\label{eq:interchange-gram}
 \left(1-\frac{k(v,w)}p\right) I_2+
   \frac{k(v,w)}p\begin{pmatrix}0&1\\1&0\end{pmatrix}.
\end{equation}
Let $\mathcal D$ be the quotient by its radical and denote the two
resulting vectors by $x_v,y_v$. Its dimension $D$ satisfies
$2\le D=O_p(n_0)$. Indeed, each summand of $k$ factors through
the $n_0$ inner words of one coordinate, and the constant matrix on
labels has rank one. Each $(x_v,y_v)$ is an orthonormal pair for the
induced form. Put $a=D-2$ and
$\mathcal D_v=\langle x_v,y_v\rangle^\perp$.

The parameter $D$ measures the cost of a general change of layout.
The two levels of the label construction brought the rank fraction
down to $O_p(q^{-p^2})$ while keeping $D$ proportional to the number
of \emph{inner} words. These rational address changes specialize
modulo $b^f$. The prime $p$ continues to govern only the scalar data.

\subsection{A shared application of a sum of swaps}

Write the $m$ pairs of main fields in physical order as
$H_1,\ldots,H_m,N_1,\ldots,N_m$, and place $a=D-2$ reserved fields
before and $a$ after them. Every field has width $f$ and ranges over
all $b^f$ values. The reserves index occupied data, so the algorithm
must restore these fields. Unmentioned spectator indices keep their
places. At each label there are $R$ independent rows, each a stream
over a fixed finite commutative $\mathbb F_p$-algebra. For
$Q_j$ the exchange of $H_j$ and $N_j$, one stage acts by
\begin{equation}\label{eq:interchange-pencil}
                   A_0+\sum_{j=1}^m A_j Q_j ,
\end{equation}
where the $A_j$ are fixed $R\times R$ scalar matrices independent of
the label. The number of calls per input stream, averaged over all
$RN$ inputs, is at most
\begin{equation}\label{eq:interchange-stage-cost}
            O\!\left(D(1+\lambda)
                  \left(1+\frac{m r_*}{N}\right)\right).
\end{equation}
All other work costs linear time in the input volume, with a constant
allowed to depend on the finite stage.

To prove this, choose a basis of $\mathcal D_v$ once for each $v$. For query
$j$ the source address map $S_{v,j}$ uses this basis for the early
reserved fields and $x_v,y_v$ for $H_j,N_j$. These $D$ columns
express the corresponding part of the address in $\mathcal D$.
The late reserved fields and all spectators retain their values. The
destination map $T_{v,j}$ uses the same basis of $\mathcal D_v$ and the
reversed columns $y_v,x_v$. To specify the action of a scalar matrix
$M$ on labels, write $X_v$ for the function giving the vector of $R$ input symbols at
each address of source $v$. For an address bijection $U$ write
$(U_*X)(d)=X(U^{-1}d)$. 
Map all source addresses into common coordinates $z$, apply the
scalar transfers there, then use the inverse destination maps to return
to physical addresses. With transfer matrix $M$ on labels this gives
the source-to-destination block $M_{wv}(T_{w,j}^{-1}S_{v,j})_*$. At a
fixed common address the output is consequently
\begin{equation}\label{eq:interchange-label-map}
       Y_w(T_{w,j}^{-1}z)=\sum_v M_{wv}\,X_v(S_{v,j}^{-1}z).
\end{equation}
For $M=I$ only label $w$ contributes to itself, and
$T_{w,j}^{-1}S_{w,j}=Q_j$. Equation~\eqref{eq:interchange-label-map}
is linear in $M$. Thus we can use the identity
\[
         I=P_*-\sum_{(L,K)}\mathbf 1_K\mathbf 1_L^{\mathsf T},
\]
where $(L,K)$ ranges over the disjoint directed rectangles of
Lemma~\ref{lem:interchange-labels}, with sources in $L$ and
destinations in $K$. The two terms may process addresses in different
physical orders. At every addition or cancellation they must identify
the same $z$.

\paragraph{Preparing one layout for all queries at an endpoint.}
A rectangle may choose its own physical order for the common
addresses. We will change basis on the reserves once at each of its
endpoints. After that preparation every query can reach the common
layout using only shifts controlled by earlier fields.

Fix a rectangle of Lemma~\ref{lem:interchange-labels} and list its
sources with orientation $(x_v,y_v)$ and its destinations with
orientation $(y_w,x_w)$. Write the resulting pairs as $(u_i,z_i)$.
By \eqref{eq:interchange-gram} every mutual oriented Gram block is
$I_2$. Choose a reference pair $(u_0,z_0)$ and set
\[
    E=\operatorname{span}\{u_i-u_0,z_i-z_0:i\},\qquad
    U=E\oplus\langle u_0,z_0\rangle .
\]
Here $E$ lies in the orthogonal complement of every oriented pair.
Thus the indicated sum is direct and $U$ is the span of all oriented
vectors. Choose a complement $C$ to $U$ in $\mathcal D$ and fix bases
of $C$ and $E$.

At endpoint $i$, lift each chosen basis vector of $C$ to the
orthogonal complement of $(u_i,z_i)$ by subtracting its projection
onto this pair. The lifts together with the basis of $E$ form a basis
of that complement. Change the early reserved fields to coordinates
in this basis. Exchange the $E$ coordinates with a prescribed equal
number of late spare fields. This change of address on the $2a$
reserved fields costs $O(D)$ calls and is independent of $j$. The
early positions now include the lifted $C$ coordinates, and the late
positions include the $E$ coordinates. For every $j$ the remaining
change to a common rectangle layout is lower triangular. In fact each
lift differs from its chosen $C$ vector by an element of $U$, and
$u_i-u_0,z_i-z_0$ lie in $E$. In physical order the common layout
therefore has $C$ coordinates early, the $u_0,z_0$ coordinates at
$H_j,N_j$, and $E$ coordinates late. All displaced spare coordinates
have the same positions at every endpoint. The forward lower
triangular change at sources and the inverse at destinations both
cost linear time. Figure~\ref{fig:rectangle-layout} displays the blocks
of this change in physical order.

\begin{figure}[!htb]
\centering
\begin{tikzpicture}[x=1cm,y=1cm,font=\small,
 stage/.style={draw=sea,fill=sea!4,rounded corners=2pt,
  text width=4.0cm,align=center,inner sep=5pt},
 head/.style={font=\footnotesize,align=center},
 block/.style={minimum width=1.22cm,minimum height=.60cm,inner sep=1pt,
   draw=ink!30,font=\small},
 idblock/.style={block,draw=sea,fill=sea!9},
 depblock/.style={block,draw=rust!75,fill=rust!8},
 zeroblock/.style={block,fill=black!2,text=ink!55}]
\node[font=\small\bfseries] at (-3.37,1.72) {Once per endpoint};
\node[stage] (lift) at (-3.37,.60)
  {Put the early reserve in a basis of lifted $C$ vectors and $E$ vectors.};
\node[stage] (late) at (-3.37,-1.20)
  {Move $E$ to the late reserve by exchanging equal-width fields.};
\draw[-{Stealth[length=1.6mm]},draw=sea] (lift)--(late);

\node[font=\small\bfseries] at (2.05,1.72) {For each query $j$};
\node[head] at (1.40,.85) {lifts of $C$};
\node[head] at (2.62,.85) {$u_i,z_i$};
\node[head] at (3.84,.85) {$E$};
\node[head,anchor=east] at (.66,.24) {$C$};
\node[head,anchor=east] at (.66,-.36) {$u_0,z_0$};
\node[head,anchor=east] at (.66,-.96) {$E$};
\foreach \x/\y in {1.40/.24,2.62/-.36,3.84/-.96}
  \node[idblock] at (\x,\y) {$I$};
\foreach \x/\y in {2.62/.24,3.84/.24,3.84/-.36}
  \node[zeroblock] at (\x,\y) {$0$};
\foreach \x/\y in {1.40/-.36,1.40/-.96,2.62/-.96}
  \node[depblock] at (\x,\y) {$*$};
\node[head] at (2.35,-1.63) {Early reserve, main pair, late reserve\\in physical order on both axes};
\end{tikzpicture}
\caption{The address change at one endpoint of a directed rectangle.
The left-hand preparation is shared by all queries. On the right,
columns are the prepared endpoint coordinates, and rows are their
coordinates in the common rectangle layout. Each possible
off-diagonal block ($*$) adds a contribution from an earlier field
to a later one, so the remaining change and its inverse require only
linear-time tape scans. Unshown fields retain their positions during
this per-query change.}
\label{fig:rectangle-layout}
\end{figure}



\paragraph{Delivery by rectangles.}
Write a directed rectangle as $(L,K)$, with sources $L$ and
destinations $K$. To implement $M=I$, evaluate
\eqref{eq:interchange-label-map} on $P_*$ and subtract its value on
each off-diagonal rectangle. For a given rectangle the quantity to
send to every $w\in K$ at a common address $z$ is the same vector
\[
             X_{L,j}(z)=\sum_{v\in L} X_v(S_{v,j}^{-1}z).
\]
Prepare the endpoint changes on the reserved fields once for every
incidence, using copies at the sources. For each $j$ the lower
triangular operations then put the summands of $X_{L,j}$ in the same
physical order. Add them in a scan, retain a copy of their sum, and
send its negative to the destinations of the rectangle. Apply $A_j$
pointwise, and at each destination accumulate the resulting terms
over $j$ after the inverse lower triangular operations. Undo the
change on the reserved fields only after this accumulation.


\paragraph{Delivery at low rank.}
For the $P_*$ contribution choose $r_*$ independent rows $I$ and
$r_*$ independent columns $J$. Over $\mathbb F_p$ we have
\[
              P_*=(P_*)_{:J}\,((P_*)_{IJ})^{-1}\,(P_*)_{I:}.
\]
This factorization uses only $r_*$ intermediate vectors at each
common address. Gathering them directly in one fixed address order
would require reordering all the source streams for each $j$ before
the first matrix could act. Instead, for basis row $i\in I$ the
vector required by $(P_*)_{I:}$ at $z$ is
\[
          X_i(S_{i,j}^{-1}z)+
            \sum_{(L,K):\,i\in K}X_{L,j}(z).
\]
The singleton uses the source layout at $i$. Every other term is a
rectangle sum already produced in its rectangle layout. Disjointness
of the rectangles accounts for each source in the row support exactly
once. There are at most $1+\lambda$ terms, irrespective of the number
of sources in the support. Convert these terms to one fixed layout
for $z$, then add. Each conversion acts on the early and late fields
and $H_j,N_j$, hence costs $O(D)$ calls per internal row. In the
fixed layout apply $((P_*)_{IJ})^{-1}$ and $A_j$ pointwise.

The support of each column in $J$ has the same decomposition on the
destination side. Convert the corresponding intermediate vector to
its singleton layout and to each incident rectangle layout, then
distribute it by the destination lower triangular operations. These
rectangle contributions enter the same destination sums before the endpoint
changes are undone. Finally add together all destined rectangle
contributions and direct singleton contributions at each destination
in its ordinary layout (along with the $A_0$ term).

Thus the changes on reserved fields cost
$O(RND(1+\lambda))$ calls for the entire stage. For each $j$ the
$r_*$ basis rows and $r_*$ basis columns use at most
$2Rr_*(1+\lambda)$ single-stream conversions. They cost
$O(mRr_*D(1+\lambda))$ calls in total, proving
\eqref{eq:interchange-stage-cost}. The row matrices are fixed. Their applications use a finite number of tapes and table entries.


The first term in the call count prepares the endpoints once for
all $m$ queries. The second is the only term that pays a fresh
coordinate change for each query. It does so on the $r_*$ intermediate
streams at each internal row index. Since $m r_*/N$ stays bounded,
these two costs are of the same order after averaging over input
streams.

\subsection{Spectral conversion and the choice of constants}
\label{subsec:swap-spectrum}

Over odd characteristic the swaps $Q_j$ have a common spectral
coordinate system. On each two-element orbit of $Q_j$ place the sum
at the address with $H_j<N_j$ and the difference (smaller address
minus larger address) at the address with $H_j>N_j$. Leave equal
pairs as they are. Apply these changes independently over all $j$.
Each conversion in either direction requires one copied interchange
per pair and pointwise operations, using $1/2$ in the inverse. The
sum and difference coordinates have eigenvalues $1$ and $-1$,
respectively. Since each positive-width field has at least two values,
one may take a two-element orbit in every pair and independently
choose its sum or difference. Thus all $2^m$ sign patterns occur.
To generate the comparisons in linear time,
fix a prefix ending immediately before $N_j$. Its $H_j$ is known,
and its $N_j$-and-suffix fibre splits into three consecutive runs.
If the suffix range is $S$, computing the run lengths takes
$\operatorname{poly}(1+f+\log S)$ steps per prefix, absorbed by the
$b^f S$ positions in the fibre. The elementary-streams lemma provides
the nested scans and amortized counters. An arbitrarily long prefix
index needs only to be advanced, while its controlled field $H_j$
uses $O(f)$ digits.


To compute the whole product $\prod_j Q_j$, use the simultaneous repair
program of \cite[Theorem~\ext{F}{thm:repair} and
Corollary~\ext{F}{cor:swap-repair}]{oaF}. It specifies the batch width
$R$ and a fixed finite scalar extension over $\mathbb F_p$, whose
elements serve as the data alphabet throughout the recursion. At
each common frequency a swap has eigenvalue $+1$ or $-1$. The
program predicts the product of these signs on most of the batch and
recovers the exceptional rows by interpolation. Its table needs only
the sign pattern, which the scans just described generate. The result
on every one of the $R$ independent input streams is exact.

The program uses $O(\sqrt{m\log m})$ stages of
\eqref{eq:interchange-pencil}, pointwise operations, and at most $4s$
full row conversions at each label, where $s\le R/(2m^3)$. Each full
conversion uses $m$ calls. The correction thus costs at most $4ms/R$
calls per input stream. The time spent in the stages is controlled by
\eqref{eq:interchange-stage-cost}.

Choose $1/2<\gamma<\alpha$. Take $q$ to be the smallest power of two
at least $m^{1/p^2}$. Then $m r_*/N=O_p(1)$,
$D=O_p(m^{(2p-1)/p^2})$, and
\[
       D(1+\lambda)\le m^{(3p-1)/p^2+o_p(1)}
          \qquad (m\longrightarrow\infty).
\]
Choose $p$ sufficiently large that $1/2+(3p-1)/p^2<\gamma$.
For sufficiently large fixed $m$, the complete computation of the
main swaps on $W=RN$ independent streams uses at most $Wm^\gamma$
calls. This follows by multiplying \eqref{eq:interchange-stage-cost}
by the number of stages and adding $4ms/R$ per stream for repair.
Also choose $m$ so that
\begin{equation}\label{eq:interchange-contraction}
                         2m^\gamma+2a<m^\alpha .
\end{equation}
All matrices, bases, row widths, and endpoint instructions have now
been fixed. Choose a compatible radix prime $b>W$, with $W>2$.
All linear operations on data, including scalar extension arithmetic,
are fixed finite tables. Recursive calls move symbols of this very
same alphabet, regardless of the depth of recursion.


\subsection{Running the recursion}
\begin{proof}[Proof of Theorem~\ref{thm:interchange}]
First consider two width-$e$ fields in radix $b$, with arbitrary
spectator ranges. Set $f=\lfloor(e-1)/(m+2a)\rfloor$. If $f=0$ the
width is bounded, and individual digit moves cost linear time. For
$f\ge1$, partition each width-$e$ field, from the front, into one
digit, $a$ fields of width $f$, the $m$ main fields of width $f$,
another $a$ fields of width $f$, and a remainder of fewer than $m+2a$
digits. The first group of $a$ fields in the earlier chunk supplies
the early reserve. The second group of $a$ fields in the later chunk
supplies the late reserve.

Split on the first digit of the earlier chunk. Group its $b$ values
into batches of $W$, padding the last batch, and retain the batch
number as a spectator. The result consists of $W$ aligned streams,
each of volume $V'\lceil b/W\rceil/b\le2V'/W$, where $V'$ is the
volume at this recursion node. Splitting, padding, and subsequent
merging cost $O_\alpha(V')$ independently of the prefix range.
Run the finite main-swap computation, implementing each call by this
same general interchange algorithm on a single stream. The complete
operation applies the main swaps to each row separately and preserves
every other address, including both occupied reserves. Discard the
padding and merge. It remains to interchange $2a$ corresponding
pairs of width-$f$ fields and boundedly many pairs of individual
digits. Perform the former recursively on the full stream and the
latter directly.

Let $T(e)$ bound time per symbol of the input stream, uniformly over
the three spectator ranges. The construction gives
\[
            T(e)\le C_\alpha+(2m^\gamma+2a)T(f).
\]
Indeed the main computation makes at most $Wm^\gamma$ calls on
streams of relative volume at most $2/W$. The other $2a$ calls
each have relative volume one. Since $f\le e/m$,
\eqref{eq:interchange-contraction} and induction give
$T(e)=O_\alpha(e^\alpha)$.

For the tape implementation use the depth-first parking schedule of
\cite[Proposition~11]{Tape}. A main call
parks the node's other live streams on stack tapes at their current
ends. Returning to the node restores them from those ends. Clearing
work areas, moving heads, and copying parked streams cost $O_\alpha(V')$
at the node, including all its finitely many descents. The child
descriptor collapses the spectator ranges to the prefix, gap, and
suffix of its two fields. Its calculation and the local control stack
use polynomial time in $1+\log V'$, absorbed by $O_\alpha(V')$.
Space at an active node is $O_\alpha(V')$. A main descent shrinks
volume by a factor at most $2/W<1$. During each of the $2a$
full-volume descents the main streams have already been merged and
their workspace cleared. Pass the active stream to the child and
park only a descriptor and return instruction. Consequently the
parked arrays sum geometrically along any path. The depth is
$O(1+\log(e+1))$ and the parked descriptors have total size
$O_\alpha((1+\log(e+1))(1+\log V'))$ at a root of volume $V'$,
which is $O_\alpha(V')$. Work tapes and their stack areas are reused,
so these bounds also count visited cells.

Finally, for binary chunks take the least $e$ with $b^e\ge 2^u$.
Pad each of the two numerical ranges from $2^u$ to $b^e$ by lexicographic
scans as in \cite[Lemma~12]{Tape}. The volume
increases by a factor less than $b^2$, and the full swap preserves
the set of genuine addresses. Scanning out those addresses returns
the required stream. This costs linear additional time and space,
and $e=O(u)$. Bits are represented by the distinct elements $0,1$ of
the data algebra. A fixed larger input alphabet can be encoded in a
fixed number of bit streams. This proves the theorem.
\end{proof}
